\section{Erd\H{o}s Problem 623 --- ROUND-2}

\subsection*{1) ROUND-2 OBJECTIVE}

\textbf{Path pursued: (C) obstruction / gap-isolation.}

Round-1 identified the main obstacle in the naive recursion: when a new point $y_n$ is chosen, the values
$f(B\cup\{y_n\})$ create uncontrolled future constraints. In this round I isolate the exact scope of that
obstruction. I will prove that \emph{every} such set-mapping on $\aleph_\omega$ has independent sets of
every finite size. Hence any negative answer to Problem 623 must be a genuinely infinitary failure of
compactness, not a finite obstruction.

\subsection*{2) Round-1 FOUNDATION USED}

I use the following Round-1 results as fixed background:

\begin{itemize}
\item the formal restatement of the problem;
\item the Round-1 observation that a direct stage-by-stage recursion toward an infinite independent set
fails because values of the form $f(B\cup\{y_n\})$ are not locally controllable;
\item the Round-1 reduction from strong closure-free principles to a positive answer, only as contextual
background (not in the proof below).
\end{itemize}

\subsection*{3) NEW INSIGHT / TOOL (ROUND-2)}

For a fixed finite target size $m$, I encode \emph{all} values of $f$ on subsets of an $m$-set into a
single finite-valued set-mapping
\[
\Phi_m:[X]^m \to [X]^{<\omega},\qquad
\Phi_m(A)=\{f(B):B\subseteq A\}.
\]
Kuratowski's free set theorem can then be applied to $\Phi_m$.
This collapses all lower arities simultaneously and yields an $(m+1)$-element set $Z$ with the property
that every proper subset of $Z$ already maps outside $Z$. Any $m$-subset of $Z$ is therefore independent
for the original map $f$.

\subsection*{4) ATTACK PLAN (ROUND-2)}

The exact gap left after Round-1 was this: we had no proof that finite independent sets exist at all sizes,
so it was unclear whether the obstruction was already visible in finite combinatorics, or only at the
infinite level.

The precise claim to prove now is:
\[
\forall m\ge 1\ \exists Y\in [X]^m\ \forall B\subseteq Y,\qquad f(B)\notin Y.
\]

If this succeeds, then the only remaining gap is the passage from ``independent sets of every finite
size'' to ``one infinite independent set,'' i.e. a pure compactness problem.

\subsection*{5) WORK (ROUND-2)}

\paragraph{Theorem.}
Let $X$ be a set of cardinality $\aleph_\omega$, and let
$f:[X]^{<\omega}\to X$ satisfy $f(A)\notin A$ for every finite $A\subseteq X$.
Then for every positive integer $m$ there exists an independent set $Y\subseteq X$ of size $m$.

\paragraph{Proof.}
Fix $m\ge 1$.
Define
\[
\Phi_m:[X]^m\to [X]^{<\omega},\qquad
\Phi_m(A):=\{f(B):B\subseteq A\}.
\]
Since $A$ is finite, $\Phi_m(A)$ is finite.

Now apply Kuratowski's free set theorem with $n=m$ and $\lambda=\aleph_0$:
because $|X|=\aleph_\omega\ge \aleph_m$, every map from $[X]^m$ to finite subsets of $X$
has a free set of size $m+1$.
Hence there exists $Z\in [X]^{m+1}$ such that
\[
\forall A\in [Z]^m\qquad \Phi_m(A)\cap Z\subseteq A.
\tag{$*$}
\]

I claim that every proper subset $C\subsetneq Z$ satisfies
\[
f(C)\notin Z.
\tag{$**$}
\]

Indeed, there are two cases.

\smallskip

\emph{Case 1: $|C|=m$.}
Then $C\in [Z]^m$, so by $(*)$ we have $\Phi_m(C)\cap Z\subseteq C$.
Since $f(C)\in \Phi_m(C)$, this implies:
if $f(C)\in Z$, then $f(C)\in C$.
But the standing hypothesis on $f$ gives $f(C)\notin C$.
Therefore $f(C)\notin Z$.

\smallskip

\emph{Case 2: $|C|<m$.}
Take any $z\in Z\setminus C$.
Then $A_z:=Z\setminus\{z\}\in [Z]^m$ and $C\subseteq A_z$, hence
$f(C)\in \Phi_m(A_z)$ by definition of $\Phi_m$.
Applying $(*)$ to $A_z$, we get
\[
\Phi_m(A_z)\cap Z\subseteq A_z=Z\setminus\{z\},
\]
so $z\notin \Phi_m(A_z)$.
Since $f(C)\in \Phi_m(A_z)$, it follows that $f(C)\neq z$.

Because $z\in Z\setminus C$ was arbitrary, we conclude that
$f(C)\notin Z\setminus C$.
Again, the hypothesis $f(C)\notin C$ shows $f(C)\notin Z$.

\smallskip

Thus $(**)$ holds for every proper subset $C\subsetneq Z$.

Now choose any $Y\in [Z]^m$.
Let $B\subseteq Y$ be arbitrary.
Because $|Z|=m+1$ and $|Y|=m$, every subset $B\subseteq Y$ is a proper subset of $Z$.
Hence by $(**)$,
\[
f(B)\notin Z,
\]
and therefore in particular $f(B)\notin Y$.
So $Y$ is independent.

This proves the theorem. \hfill $\square$

\paragraph{Corollary.}
If Problem 623 has a negative answer, then the obstruction is purely infinitary:
for every $m<\omega$ there are independent sets of size $m$, but there is no infinite independent set.

\subsection*{6) ADVERSARIAL VERIFICATION}

\begin{itemize}
\item \emph{Potential issue: does the proof handle the case $B=Y$?}
Yes. Since $Y$ has size $m$ and $Z$ has size $m+1$, the set $Y$ is a proper subset of $Z$,
so $(**)$ applies directly and gives $f(Y)\notin Z$, hence $f(Y)\notin Y$.

\item \emph{Potential issue: did Kuratowski's theorem apply to the correct arity?}
Yes. The map $\Phi_m$ is defined on $m$-element subsets, and the theorem gives a free set of
size $m+1$ because $|X|=\aleph_\omega\ge \aleph_m$.

\item \emph{Potential issue: could $\Phi_m(A)$ fail to be finite?}
No. Each $A$ has exactly $2^m$ subsets, so $\Phi_m(A)$ has size at most $2^m$.

\item \emph{Interaction with Round-1.}
Round-1 isolated a recursion gap for building an \emph{infinite} independent set.
The theorem above shows that this gap cannot be repaired merely by improving finite-stage bookkeeping:
all finite target sizes are always attainable. The unresolved difficulty is exactly the infinite compactness step.
\end{itemize}

\subsection*{7) FINAL}

\textbf{UNRESOLVED (BUT STRICTLY ADVANCED).}

The strict advance over Round-1 is the unconditional theorem that
\emph{independent sets of every finite size always exist}.
So any counterexample to Problem 623 must have arbitrarily large finite independent sets and yet no infinite one.
That sharply localizes the remaining obstruction to a genuinely infinitary compactness phenomenon.

\subsection*{8) COMPLETION ESTIMATE (MANDATORY)}

COMPLETION: 40\%

\subsection*{9) REFERENCES}

\begin{itemize}
\item T. F. Bloom, \emph{Erd\H{o}s Problem \#623}, \texttt{erdosproblems.com/623}, accessed 2026-03-07.
\item K. Kuratowski, \emph{Sur une caract\'erisation des alephs}, Fund. Math. \textbf{38} (1951), 14--17.
\end{itemize}

\subsection*{10) OUTPUT FORMAT}

This section is written in drop-in \LaTeX\ form and starts directly from \verb|\section{}|, with no preamble.

\bigskip

